% 编译：xelatex math-foundations.tex（需要 ctex 宏包，编译两遍以解析引用）
\documentclass[11pt,a4paper]{ctexart}
\usepackage[margin=2.5cm]{geometry}
\usepackage{amsmath,amssymb}
\usepackage{booktabs,tabularx}
\usepackage{enumitem}
\usepackage[hidelinks]{hyperref}
\usepackage{xurl}

\newcommand{\E}{\mathbb{E}}
\DeclareMathOperator{\Var}{Var}
\DeclareMathOperator{\Cov}{Cov}
\DeclareMathOperator{\sgn}{sgn}
\DeclareMathOperator{\ReLU}{ReLU}
\DeclareMathOperator{\round}{round}

\newenvironment{caution}{\par\medskip\noindent\textbf{注意.}\ }{\par\medskip}

\title{跨层误差传播模型：数学基础}
\date{2026年11月12日}

\begin{document}
\maketitle

\begin{abstract}
本文给出跨层误差传播模型的数学内容，按三个层间传播算子组织：$\mathcal{P}_{12}$（器件$\to$电路）、$\mathcal{P}_{23}$（电路$\to$框架）、$\mathcal{P}_{34}$（框架$\to$应用）。在此基础上给出算子组合的误差界与接口充分性的定义（对应研究问题 R1、R3），以及逐边界的掩蔽与增益度量（对应 R4）。目标、范围与数制约定见 \texttt{overview.tex}。
\end{abstract}

\tableofcontents

\section{记号与参考模型}

\begin{table}[h]
\centering
\small
\begin{tabularx}{\textwidth}{@{}lX@{}}
\toprule
记号 & 含义 \\
\midrule
$a, b$ & 乘法器操作数：$a$ 为有符号 INT8 权重，$b$ 为无符号 8 位激活 \\
$m(a,b)=ab$，$\tilde m(a,b)$ & 精确乘法、近似乘法 \\
$e(a,b)$ & 近似单元的设计误差 $\tilde m(a,b) - ab$（电路层误差源 $S_2$） \\
$\pi_k$ & 第 $k$ 个故障位置在一次运算中出错的概率（器件层误差源 $S_1$） \\
$p_x$ & 某网络层激活的分布 \\
$N$ & 一个点积的累加长度 \\
$\Delta_i$ & 接口 $E_i$ 上的误差（相对参考模型） \\
$R_i = \rho_i(\Delta_i)$ & 接口 $E_i$ 上误差的表示（摘要） \\
$\Delta^{(l)}$ & 第 $l$ 个网络层输出的误差向量 \\
$J_l$ & 第 $l$ 个网络层对其输入的雅可比矩阵 \\
$z$，$c$ & 输出 logits；参考模型预测的类别 \\
$\Phi, \phi$ & 标准正态分布函数与密度 \\
\bottomrule
\end{tabularx}
\caption{记号}
\end{table}

\paragraph{参考模型.}
所有误差都相对于\textbf{精确的 INT8 量化模型}定义（数制约定见总纲），而不是相对于浮点模型。因此量化本身的误差不计入，除非作为框架层误差源 $S_3$ 被显式引入（如累加器截断）。

\section{$\mathcal{P}_{12}$：器件层 $\to$ 电路层}

\subsection{故障位置与单故障假设}
运算单元中的故障位置 $k$ 包括：操作数寄存器的各位、乘积寄存器的各位、累加器的各位、单元内部的逻辑节点。设每次运算中位置 $k$ 出错的概率为 $\pi_k$，且 $\sum_k \pi_k \ll 1$，则可假设一次运算至多一个故障。

\subsection{存储位翻转}
\paragraph{激活位翻转.} 无符号激活 $b$ 的第 $j$ 位翻转：
\begin{equation}
\Delta b = (1 - 2b_j)\,2^j,\qquad
\Delta y = \tilde m(a, b + \Delta b) - \tilde m(a, b).
\end{equation}
对精确乘法器，$\Delta y = a\,\Delta b$：位错误被\textbf{另一个操作数放大}，误差依赖数据。

\paragraph{权重位翻转.} 有符号补码权重 $a$ 的第 $j$ 位翻转，位权 $w_j = 2^j$（$j<7$），$w_7 = -2^7$：
\begin{equation}
\Delta a = (1 - 2a_j)\,w_j .
\end{equation}

\paragraph{乘积或累加器位翻转.} 直接得到 $\Delta = (1 - 2y_j)\,2^j$。32 位累加器的高位翻转会产生极大误差，属于“罕见大误差”。

\subsection{逻辑节点错误与逻辑掩蔽}
内部节点 $k$ 出错后，输出误差由门级网表决定：
\begin{equation}
\Delta y = g_k(a,b) - \tilde m(a,b),
\end{equation}
其中 $g_k$ 是节点 $k$ 出错时的电路函数。\textbf{逻辑掩蔽概率}
\begin{equation}
\mu^{\text{mask}}_k = P(\Delta y = 0 \mid \text{节点 } k \text{ 出错})
\end{equation}
通过对网表穷举 65536 种输入并逐节点注入求得。故障应注入在\textbf{近似单元本身的网表}上，而不是精确单元上。

\subsection{时序违例}
若寄存器在违例时捕获旧值，第 $k$ 位的误差为 $\delta_k = y_k^{\text{prev}} - y_k \in \{-1,0,1\}$，只在本次输出位与上次不同时可能非零，因此依赖输入序列。时序违例是可选的物理误差源，其出错概率模型见附录~\ref{app:timing}。

\subsection{E2 的构成与摘要}
在单故障假设下，给定操作数 $(a,b)$，单元输出误差的分布为
\begin{equation}
P(\Delta y = d \mid a,b) = \Bigl(1 - \sum_k \pi_k\Bigr)\,\mathbf{1}[d = e(a,b)] + \sum_k \pi_k\, \mathbf{1}\bigl[d = e(a,b) + \Delta_k(a,b)\bigr],
\end{equation}
其中 $\Delta_k(a,b)$ 是故障 $k$ 引起的额外误差。E2 的摘要分两部分：
\begin{itemize}[leftmargin=*]
  \item \textbf{小误差分量}：以权重为条件、对激活分布和故障分布取期望的矩
  \begin{equation}
  \mu(a) = \E_{b \sim p_x,\,k}[\Delta y],\qquad \sigma^2(a) = \Var_{b \sim p_x,\,k}[\Delta y]
  \end{equation}
  （只计入幅度不超过阈值 $\tau$ 的误差）；
  \item \textbf{大误差分量}：$\lambda(a) = P(|\Delta y| > \tau)$ 及其幅度分布 $F_{\text{big}}$。
\end{itemize}
阈值 $\tau$ 是模型参数，按“小误差分量的正态近似是否仍然成立”选取。

\begin{caution}
\begin{itemize}[leftmargin=*]
  \item 推理时权重固定，必须以权重为条件，不能对权重分布再平均。
  \item 有符号乘法用符号—幅值包装时，$e_{\text{s}}(a,b) = \sgn(a)\,e_{\text{u}}(|a|,b)$；一个点积的偏置 $\sum_i \sgn(w_i)\,\mu_{\text{u}}(|w_i|)$ 在正负权重平衡时大幅抵消。
  \item 存储位翻转的误差随另一操作数放大，因此即使 $\pi_k$ 与数据无关，E2 也依赖数据。
\end{itemize}
\end{caution}

\section{$\mathcal{P}_{23}$：电路层 $\to$ 框架层}

\subsection{点积}
$y = \sum_{i=1}^{N} w_i x_i$，$\Delta y = \sum_i \varepsilon_i$。整数累加精确时
\begin{equation}
\E[\Delta y] = \sum_{i} \mu(w_i),\qquad
\Var[\Delta y] = \sum_{i} \sigma^2(w_i) + \sum_{i \neq j} \Cov(\varepsilon_i, \varepsilon_j).
\end{equation}
\paragraph{偏置主导的判据.} 独立同分布时 $R = \sqrt N\,|\mu|/\sigma$；当 $N > N^{*} = \sigma^2/\mu^2$ 时偏置主导。

\paragraph{协方差的三个来源.}
\begin{enumerate}[leftmargin=*]
  \item 点积内部：相邻激活相关，$\Cov(\varepsilon_i,\varepsilon_j) \ne 0$；
  \item 输出之间：卷积的不同输出位置共用权重，$\Cov(\Delta y_p, \Delta y_q) = \sum_{i,j}\Cov\bigl(e(w_i,x_{i,p}), e(w_j,x_{j,q})\bigr)$；
  \item 传播误差与新增误差之间：见第~\ref{sec:recursion} 节。
\end{enumerate}

\subsection{框架层误差源 $S_3$：累加器截断与饱和}
\paragraph{截断低位.} 部分和向下截断低 $t$ 位时，若被截掉的低位近似均匀分布，
\begin{equation}
\E[\Delta] = -\frac{2^t - 1}{2},\qquad \Var[\Delta] = \frac{2^{2t} - 1}{12}.
\end{equation}
截断是\textbf{有偏的}，与舍入不同。
\paragraph{饱和.} $B$ 位累加器的饱和概率 $P_{\text{sat}} \approx P(|y + \Delta y| > 2^{B-1})$；饱和既截断误差，也引入大误差。A2Q \cite{a2q} 给出的权重 $L_1$ 范数界可用于判断是否可能溢出。

\subsection{重量化舍入：掩蔽}
重量化 $q = \round(s\cdot\text{acc} + \beta)$。设缩放后误差为 $\delta$，写成 $\delta = n + \varphi$，$n = \lfloor\delta\rfloor$，$\varphi\in[0,1)$。若舍入前数值的小数部分近似均匀且与 $\delta$ 独立，则输出误差 $r = \round(z+\delta) - \round(z) \in \{n, n+1\}$，且
\begin{equation}
P(r = n+1) = \varphi,\qquad \E[r\mid\delta] = \delta,\qquad \Var[r\mid\delta] = \varphi(1-\varphi).
\end{equation}
当 $|\delta| < 1$ 时，$P(r = 0) = 1 - |\delta|$，即小误差以概率 $1-|\delta|$ 被\textbf{掩蔽}。

\begin{caution}
舍入保持误差均值，但改变二阶矩：$\E[r^2\mid\delta] = \delta^2 + \varphi(1-\varphi)$。线性传播 $r\approx\delta$ 会低估方差。小数部分均匀的假设在缩放因子特殊时可能不成立，需实测检验。
\end{caution}

\subsection{非线性与结构}
\paragraph{一般方法.} delta 方法：$\E[g(u+\Delta)] - g(u) \approx g'(u)\mu_\Delta + \tfrac12 g''(u)(\sigma_\Delta^2 + \mu_\Delta^2)$，$\Var \approx g'(u)^2\sigma_\Delta^2$；向量形式 $\Sigma_{\text{out}} \approx J\Sigma_{\text{in}}J^{\mathsf T}$。

\paragraph{ReLU.} 给定精确输入 $u$、误差 $\Delta\sim\mathcal N(\mu,\sigma^2)$，令 $m = u+\mu$：
\begin{equation}
\E[\ReLU(u+\Delta)] = m\,\Phi(m/\sigma) + \sigma\,\phi(m/\sigma).
\end{equation}
$u$ 远小于 0 时误差被完全掩蔽。

\paragraph{推理时折叠的 BatchNorm.} 逐通道仿射，$\Delta y = (\gamma/s)\Delta x$；\textbf{不吸收偏置}，只有重新标定统计量或重训练后才会。

\paragraph{残差.} $\Delta y = \Delta F + \Delta x$，$\Var[\Delta y] = \Var[\Delta F] + \Var[\Delta x] + 2\Cov(\Delta F,\Delta x)$；误差绕过中间层的掩蔽。

\paragraph{池化.} 平均池化为线性，方差缩小程度取决于窗口内误差的相关性；最大池化在最大值与次大值差距远大于误差尺度时 $\Delta y \approx \Delta x_{i^*}$。

\paragraph{Softmax 与 LayerNorm（后续阶段）.} Softmax 的雅可比 $J = \operatorname{diag}(p) - pp^{\mathsf T}$，$J\mathbf 1 = 0$，共模偏置被掩蔽；LayerNorm 的雅可比 $\frac{\operatorname{diag}(\gamma)}{\sigma_x}\bigl(I - \frac1d\mathbf 1\mathbf 1^{\mathsf T} - \frac1d\hat x\hat x^{\mathsf T}\bigr)$ 同样去掉共模偏置。

\subsection{层间递推}
\label{sec:recursion}
记精确激活 $a_l$，第 $l$ 层的近似实现 $\tilde f_l$：
\begin{equation}
\Delta^{(l+1)} = \underbrace{f_l(a_l+\Delta^{(l)}) - f_l(a_l)}_{\approx J_l\Delta^{(l)}} + \underbrace{\tilde f_l(a_l+\Delta^{(l)}) - f_l(a_l+\Delta^{(l)})}_{\varepsilon^{(l)}},
\end{equation}
\begin{align}
m_{l+1} &\approx J_l m_l + \mu_\varepsilon^{(l)},\\
\Sigma_{l+1} &\approx J_l\Sigma_l J_l^{\mathsf T} + \Sigma_\varepsilon^{(l)} + J_l C_l + C_l^{\mathsf T}J_l^{\mathsf T},\qquad C_l = \Cov(\Delta^{(l)},\varepsilon^{(l)}).
\end{align}

\subsection{协方差的表示：三档}
完整协方差矩阵维度为（通道 $\times$ 高 $\times$ 宽）的平方，必须选择表示方式，这正是 R1 要比较的：
\begin{enumerate}[leftmargin=*]
  \item \textbf{对角}（ADF \cite{gast2018}）：只传方差。
  \item \textbf{分组概率仿射算术}：误差写成
  \begin{equation}
  \Delta = c + \sum_{s} \alpha_s\,\epsilon_s,\qquad \Cov(\Delta_p,\Delta_q) = \sum_s \alpha_{p,s}\alpha_{q,s}\Var(\epsilon_s),
  \end{equation}
  即协方差以 $A D A^{\mathsf T}$ 的分解形式保存。噪声符号 $\epsilon_s$ 按误差来源分组（如每个“网络层，输出通道”一个），贡献小的符号定期合并。线性层精确；ReLU 跨零时按 DeepZ 方式线性化并新增符号；舍入不是仿射运算，按上面的掩蔽公式处理，新增方差 $\varphi(1-\varphi)$ 记为新的独立符号。
  \item \textbf{仿真真值}：逐比特仿真得到的完整统计。
\end{enumerate}

\subsection{大误差分量的传播}
大误差（幅度超过 $\tau$）罕见但往往直接改变输出。对第 $l$ 个网络层定义\textbf{敏感度}
\begin{equation}
s_l = P(\text{输出改变} \mid \text{第 } l \text{ 层发生一次大误差}),
\end{equation}
由少量故障注入估计（对应 Vallero 等人的条件概率表）。E3 中记录每层大误差发生率 $\Lambda_l = n_l\,\bar\lambda_l$（$n_l$ 为该层运算次数）。

\begin{caution}
\begin{itemize}[leftmargin=*]
  \item 新增误差依赖已被污染的输入，用精确激活下的统计量代替是一阶近似；交叉项 $C_l$ 通常假设为 0，需用仿真检验。
  \item $J_l$ 依赖数据，实际使用的是对数据分布取期望的平均场近似。
  \item $s_l$ 会受小误差分量影响（小误差改变了 logit 间隔），“$s_l$ 与配置无关”只是近似。
\end{itemize}
\end{caution}

\section{$\mathcal{P}_{34}$：框架层 $\to$ 应用层}
本节思路与 Sakr 等人的“失配概率” \cite{sakr2017} 相同；扩展在于加入偏置、logits 间协方差与大误差分量。

\paragraph{小误差分量.} 样本 $s$ 的间隔 $M_{s,j} = z_c - z_j$，$D_j = \Delta z_j - \Delta z_c$：
\begin{equation}
P_{\text{small},s} \le \sum_{j\ne c}\Phi\!\left(\frac{\E[D_j] - M_{s,j}}{\kappa\,\sigma_{D_j}}\right),\qquad
\sigma_{D_j}^2 = \Var[\Delta z_j] + \Var[\Delta z_c] - 2\Cov(\Delta z_j,\Delta z_c),
\end{equation}
$\kappa$ 为尾部标定因子。只有 logits 误差的差起作用，共模偏置被掩蔽。

\paragraph{大误差分量.}
\begin{equation}
P_{\text{big}} = 1 - \prod_l (1 - s_l)^{\Lambda_l} \approx \sum_l \Lambda_l s_l .
\end{equation}

\paragraph{合并与质量.} 假设两分量独立，$P_{\text{flip},s} = 1 - (1-P_{\text{small},s})(1-P_{\text{big}})$；
\begin{equation}
\widehat{\Delta\text{Acc}} \approx \frac{1}{S}\sum_{s\in\mathcal C} P_{\text{flip},s},
\end{equation}
$\mathcal C$ 为参考模型分类正确的样本。与参考模型的预测一致率 $1 - \frac1S\sum_s P_{\text{flip},s}$ 不受标签噪声影响，是更干净的目标。

\section{算子组合与接口充分性（R1、R3）}
\label{sec:composition}

\subsection{定义}
记真实传播为 $T_i:\Delta_i\mapsto\Delta_{i+1}$，接口表示为 $\rho_i$，模型算子为 $\mathcal P_i: R_i \mapsto R_{i+1}$。理想情况下
\begin{equation}
\mathcal P_i\circ\rho_i = \rho_{i+1}\circ T_i .
\end{equation}
定义算子 $i$ 的\textbf{充分性缺陷}（在表示空间的度量 $d$ 下，对实际出现的误差取平均或最大）：
\begin{equation}
\eta_i = d\bigl(\mathcal P_i(\rho_i(\Delta_i)),\ \rho_{i+1}(T_i(\Delta_i))\bigr).
\end{equation}
若 $\rho_i$ 是“充分”的（$T_i$ 的输出摘要只通过 $\rho_i$ 依赖输入误差），且 $\mathcal P_i$ 准确，则 $\eta_i = 0$。

\subsection{组合误差界}
设 $\mathcal P_j$ 在度量 $d$ 下的 Lipschitz 常数为 $L_j$。由逐项拆分（望远镜求和）得
\begin{equation}
d\bigl(\mathcal P_{3}\circ\mathcal P_{2}\circ\mathcal P_{1}(R_1),\ \rho_4(T_3 T_2 T_1(\Delta_1))\bigr) \le \sum_{i}\Bigl(\prod_{j>i}L_j\Bigr)\eta_i .
\end{equation}
即：组合误差不超过各算子充分性缺陷经上层放大后之和。R3 的实验检验这个界是否紧，以及哪一项占主导。

\subsection{充分性缺陷的来源}
\begin{enumerate}[leftmargin=*]
  \item \textbf{丢失数据依赖}：用均匀分布或对权重平均后的矩代替以权重为条件的矩。
  \item \textbf{丢失相关性}：只传对角方差。
  \item \textbf{非线性}：一阶近似在误差较大或靠近 ReLU 拐点时失效。
  \item \textbf{非正态尾部}：罕见大误差被并入方差。
\end{enumerate}
R1 的实验逐一去掉上述信息，测量 $\eta_i$ 的变化。

\section{掩蔽与增益度量（R4）}
对一个边界，输入误差 $\Delta_{\text{in}}$、输出误差 $\Delta_{\text{out}}$，定义：
\begin{align}
\text{掩蔽概率}\quad & m = P(\Delta_{\text{out}} = 0 \mid \Delta_{\text{in}} \ne 0),\\
\text{能量增益}\quad & G = \E\|\Delta_{\text{out}}\|^2 \,/\, \E\|\Delta_{\text{in}}\|^2,\\
\text{偏置传递}\quad & B = \langle\E[\Delta_{\text{out}}],\,v\rangle \,/\, \langle\E[\Delta_{\text{in}}],\,u\rangle,
\end{align}
其中 $u, v$ 为选定的投影方向（如全 1 向量或对 logit 差的方向）。

\paragraph{分解.} 在组合成立、一阶近似有效时：
\begin{equation}
G_{\text{total}} = \prod_i G_i \ \Rightarrow\ \log G_{\text{total}} = \sum_i \log G_i,\qquad
1 - m_{\text{total}} = \prod_i (1 - m_i),
\end{equation}
各边界的贡献份额为 $\log G_i/\log G_{\text{total}}$。

\begin{caution}
分解只在组合成立时严格有效；实测的 $G_{\text{total}}$ 与 $\prod_i G_i$ 之差本身就是 R3 的一个度量。
\end{caution}

\section{统计与验证}
\begin{itemize}[leftmargin=*]
  \item \textbf{出错概率}：Wilson 置信区间；零次观测时 95\% 上界约 $3/n$。
  \item \textbf{精度分辨率}：$\text{SE} = \sqrt{A(1-A)/S}$；CIFAR-10 上约 0.3\%，两次独立测量 0.8\% 以内的差异无法分辨。
  \item \textbf{配对比较}：统计“正确变错误” $n_{10}$ 与“错误变正确” $n_{01}$，$\Delta\text{Acc} = (n_{10}-n_{01})/S$，McNemar 检验。
  \item \textbf{预测评价}：矩用对数比；质量用平均绝对误差与最大误差；排序用 Spearman 秩相关。
  \item \textbf{标定与验证分离}：$\kappa$、$\tau$、$s_l$ 的标定数据不得用于验证。
\end{itemize}

\section{注意点汇总}
\begin{enumerate}[leftmargin=*]
  \item 误差相对精确 INT8 模型定义，量化本身不计入。
  \item 推理时权重固定，使用以权重为条件的矩；有符号数处理方式按总纲约定。
  \item 存储位翻转的误差随另一操作数放大，E2 依赖数据。
  \item 故障注入在近似单元自身的网表上。
  \item 累加器截断有偏，舍入无偏；舍入掩蔽小误差并改变二阶矩。
  \item 推理时折叠的 BatchNorm 不吸收偏置。
  \item Softmax、argmax 掩蔽共模偏置，应传播 logits 误差的差。
  \item 协方差有三处来源；对角近似会低估池化与累加后的误差。
  \item 罕见大误差不能并入方差，单独以敏感度 $s_l$ 处理。
  \item 组合与分解只在接口充分时成立；充分性缺陷要实测。
  \item 精度测量本身有约 0.3\% 的标准误差，用配对比较。
\end{enumerate}

\appendix
\section{可选误差源：时序违例的出错概率模型}
\label{app:timing}
第 $k$ 位路径延迟 $D_k$ 超过 $T_{\text{clk}} - t_{\text{su}}$ 时违例：
\begin{equation}
\pi_k(f,V,T) = P\bigl(D_k > T_{\text{clk}} - t_{\text{su}}\bigr) \approx \Phi\!\left(\frac{\mu_k(V,T) - (T_{\text{clk}} - t_{\text{su}})}{\sigma_k(V,T)}\right),
\end{equation}
延迟随电压可用 alpha 幂律 $\mu_k \propto V/(V - V_{\text{th}})^{\alpha}$ 拟合。注意：这只是拟合形式；先进工艺可能出现温度反转；出错概率依赖输入跳变，需要一个降维特征（如输出翻转位数）作为条件；器件数据必须在做端到端验证的同一块板上测得。

\begin{thebibliography}{9}
\bibitem{gast2018} J.~Gast, S.~Roth. Lightweight probabilistic deep networks. CVPR, 2018.
\bibitem{sakr2017} C.~Sakr, Y.~Kim, N.~Shanbhag. Analytical guarantees on numerical precision of deep neural networks. ICML, 2017.
\bibitem{a2q} I.~Colbert, A.~Pappalardo, J.~Petri-Koenig. A2Q: accumulator-aware quantization with guaranteed overflow avoidance. ICCV, 2023.
\bibitem{vallero2016} A.~Vallero et al. Cross-layer system reliability assessment framework for hardware faults. ITC, 2016.
\end{thebibliography}

\end{document}
